% !TEX root = ../main.tex
% School-provided disclosure form reproduced in LaTeX so it is compiled with
% the thesis.  The wording follows the supplied DOCX; the signatory is left
% blank for handwritten signatures.

\begingroup
\newgeometry{left=2.70cm,right=2.70cm,top=3.70cm,bottom=2.35cm,headheight=0pt,headsep=0pt,footskip=0pt}
\pagestyle{empty}
\flushbottom
\fontspec{Times New Roman}
\CJKfontspec[BoldFont={Noto Serif CJK SC Bold}]{Noto Serif CJK SC}
\setlength{\parindent}{2em}
\setlength{\parskip}{0pt}

\newcommand{\AIStatementPageNumber}[1]{%
  \vfill
  \begin{center}\fontsize{10.5bp}{12bp}\selectfont —#1—\end{center}%
  \clearpage
}

% Page 1
\thispagestyle{empty}
{\fontsize{14bp}{20bp}\selectfont\noindent 附件 2\par}
\vspace{1.0cm}
{\centering\fontsize{21bp}{30bp}\selectfont 毕业设计（论文）中人工智能工具使用说明\par}
\vspace{0.75cm}
{\fontsize{16bp}{29bp}\selectfont
依据《上海交通大学关于在教育教学中使用 AI 的规范》的要求，为充分发挥“AI+HI（人工智能+人类智慧，Artificial Intelligence + Human Intelligence）”在毕业设计（论文）中的作用，特制定《毕业设计（论文）中人工智能工具使用说明》。该说明旨在支持人工智能工具在毕业设计（论文）中发挥辅助与启发功能，同时通过合理披露使用方式，确保毕业设计（论文）的完成过程中能够充分发挥“人类智慧”的作用，并符合科研诚信的要求。\par
}
\AIStatementPageNumber{1}

% Page 2
\thispagestyle{empty}
{\fontsize{11bp}{21bp}\selectfont\bfseries
\noindent 毕业设计（论文）标题：MNEME：面向持续文献理解的移动原生人工智能研究助手\par
\vspace{0.20cm}
\noindent 学生姓名：王瀚洋、赵恒、蒋睿宇、张轶凡\par
\vspace{0.20cm}
\noindent 学号：522370910183、522010910074、522370910078、522370910034\par
\vspace{0.20cm}
\noindent 指导教师：皮宜博\par
}
\vspace{0.65cm}
{\centering\fontsize{16bp}{22bp}\selectfont\bfseries 一、本文使用的 AI 工具清单\par}
\vspace{0.30cm}
{\fontsize{8.2bp}{14bp}\selectfont
\setlength{\tabcolsep}{2.4pt}
\renewcommand{\arraystretch}{1.25}
\arrayrulecolor{black}
\begin{tabular}{|>{\centering\arraybackslash}m{0.75cm}|>{\centering\arraybackslash}m{1.65cm}|>{\centering\arraybackslash}m{1.20cm}|>{\centering\arraybackslash}m{1.65cm}|m{7.65cm}|}
\hline
\textbf{序号\textsuperscript{a}} & \textbf{AI 工具名称} & \textbf{版本/型号} & \textbf{毕业设计（论文）中应用部分\textsuperscript{b}} & \centering\arraybackslash\textbf{人为判断过程\textsuperscript{c}说明} \\
\hline
1 & Anthropic Claude & Claude Opus 4.8 & B3、C1 & 输出均限定于指定论文版本和检索证据；根据固定测试用例、原文、引用标记、来源匹配及拒答结果人工核验，实验数据和结论由作者独立分析。 \\
\hline
2 & Anthropic Claude & Claude Haiku 4.5 & B3、C1 & 在相同测试集和检索配置下开展对照；人工检查结构化输出、引用来源、证据匹配和拒答行为，并依据原始延迟与成本记录分析结果。 \\
\hline
\end{tabular}
\par
}
\AIStatementPageNumber{2}

% Page 3
\thispagestyle{empty}
{\fontsize{8.2bp}{14bp}\selectfont
\setlength{\tabcolsep}{2.4pt}
\renewcommand{\arraystretch}{1.25}
\arrayrulecolor{black}
\begin{tabular}{|>{\centering\arraybackslash}m{0.75cm}|>{\centering\arraybackslash}m{1.65cm}|>{\centering\arraybackslash}m{1.20cm}|>{\centering\arraybackslash}m{1.65cm}|m{7.65cm}|}
\hline
3 & OpenAI Codex & GPT-5.6 & C1、C2、D3、E1 & 生成或修改的实验/测试脚本和功能代码均经逐段审查，并通过单元测试、集成测试、静态检查和实际运行验证；润色内容逐句核对，LaTeX 排版通过编译结果、学校格式规范及最终 PDF 检查。 \\
\hline
\end{tabular}
\par
}
\vspace{0.32cm}
{\fontsize{9.5bp}{16bp}\selectfont
\noindent 注释：\par
\noindent\textsuperscript{a} 可以根据实际情况增删行数；\par
\noindent\textsuperscript{b} 可以对照“应用部分对照表”填写对应的编号表，若超过编号表内容，请根据实际填写；\par
\noindent\textsuperscript{c} 人为判断过程指的是依据何种方式判断 AI 工具生成的内容的准确性。\par
}
\vspace{0.65cm}
{\centering\fontsize{14bp}{20bp}\selectfont\bfseries 毕业设计（论文）中应用部分对照表\par}
\vspace{0.28cm}
{\fontsize{8.9bp}{14bp}\selectfont
\setlength{\tabcolsep}{4pt}
\renewcommand{\arraystretch}{1.25}
\arrayrulecolor{black}
\begin{tabular}{|m{3.4cm}|m{10.0cm}|}
\hline
A.研究准备阶段 & A1.选题与方向设定；A2.文献检索与筛选 \\
\hline
B.研究设计阶段 & B1.研究思路设计；B2.方法论构建；B3.研究工具开发 \\
\hline
C.研究实施阶段 & C1.实验/数据分析；C2.代码编写与调试；C3.数学建模与计算 \\
\hline
D.论文写作阶段 & D1.大纲与逻辑优化；D2.论文正文写作；D3.语言与语法润色 \\
\hline
E.其他辅助环节 & E1.参考文献与格式；E2.图表/多媒体生成；E3.学术反馈与模拟答辩 \\
\hline
\end{tabular}
}
\AIStatementPageNumber{3}

% Page 4
\thispagestyle{empty}
{\fontsize{14bp}{22bp}\selectfont\bfseries\noindent 二、使用说明\par}
\vspace{0.15cm}
{\fontsize{12bp}{22bp}\selectfont
（针对对应序号，简述人工智能工具如何辅助和启发研究，例如：“1. A2：DeepSeek 协助综述部分框架搭建，关键观点均经文献核实；2. C2：GitHub Copilot 用于 Python 数据清洗代码的初步生成，人工调整后实现功能。”）\par
}
\vspace{0.28cm}
{\fontsize{10.5bp}{21bp}\selectfont
\noindent 1. B3、C1：Claude Opus 4.8 作为 MNEME 的旗舰模型，用于结构化论文摘要、基于证据的单篇论文问答及固定测试集上的系统评价。生成结果均与指定论文版本、检索片段及预设指标核对；异常结果结合原文人工复核，实验数据及结论由作者独立分析。\par
\vspace{0.38cm}
\noindent 2. B3、C1：Claude Haiku 4.5 用于问答与论文摘要任务的对照实验。作者在相同测试集和检索配置下检查其结构化输出、引用来源、证据匹配与拒答行为，并依据原始延迟和成本记录分析实验结果。\par
\vspace{0.38cm}
\noindent 3. C1、C2、D3、E1：OpenAI Codex（GPT-5.6）用于辅助编写实验/测试脚本、部分功能代码及代码调试，并用于论文语言润色和 LaTeX 排版。代码由作者逐段审查并通过单元测试、集成测试、静态检查和实际运行验证；润色内容逐句核对，LaTeX 修改通过编译结果、学校格式规范及最终 PDF 页面检查。\par
}
\AIStatementPageNumber{4}

% Page 5
\thispagestyle{empty}
\vspace*{0.45cm}
{\centering\fontsize{18bp}{28bp}\selectfont\bfseries 学术诚信声明\par}
\vspace{0.55cm}
{\fontsize{13bp}{25bp}\selectfont\bfseries
本人承诺：论文核心观点、实验数据及结论均独立完成，AI 工具仅用于辅助性工作，已明确说明所有生成内容并在文中对应位置进行脚注标注。\par
}
\vspace{0.95cm}
{\fontsize{12bp}{22bp}\selectfont\bfseries
\hspace*{8.7cm}承诺人：\par
\vspace{0.35cm}
\hspace*{6.7cm}日期：2026 年 8 月 10 日\par
}
\vfill
{\fontsize{9.5bp}{16bp}\selectfont
\hrule
\vspace{0.12cm}
\noindent 上海交通大学党政办公室\hfill 主动公开\hfill 2025 年 7 月 9 日印发\par
\vspace{0.12cm}
\hrule
}
\vspace{0.35cm}
\begin{center}\fontsize{10.5bp}{12bp}\selectfont —5—\end{center}

\restoregeometry
\endgroup
